
\documentclass[final]{beamer}

\usepackage[scale=1.24]{beamerposter} % Use the beamerposter package for laying out the poster
\usepackage{bbm}
\usepackage{graphicx}
\usepackage{textpos}
\usepackage{xcolor}
\graphicspath{{./figs/}}
\usetheme{confposter} % Use the confposter theme supplied with this template

\setbeamercolor{block title}{fg=illiniblue,bg=white} % Colors of the block titles
\setbeamercolor{block body}{fg=black,bg=white} % Colors of the body of blocks
\setbeamercolor{block alerted title}{fg=white,bg=illiniorange} % Colors of the highlighted block titles
\setbeamercolor{block alerted body}{fg=black,bg=heritageorange} % Colors of the body of highlighted blocks
% Many more colors are available for use in beamerthemeconfposter.sty

%-----------------------------------------------------------
% Define the column widths and overall poster size
% To set effective sepwid, onecolwid and twocolwid values, first choose how many columns you want and how much separation you want between columns
% In this template, the separation width chosen is 0.024 of the paper width and a 4-column layout
% onecolwid should therefore be (1-(# of columns+1)*sepwid)/# of columns e.g. (1-(4+1)*0.024)/4 = 0.22
% Set twocolwid to be (2*onecolwid)+sepwid = 0.464
% Set threecolwid to be (3*onecolwid)+2*sepwid = 0.708

\newlength{\sepwid}
\newlength{\onecolwid}
\newlength{\twocolwid}
\newlength{\threecolwid}
\newlength{\fourcolwid}
\setlength{\paperwidth}{48in} % A0 width: 46.8in
\setlength{\paperheight}{36in} % A0 height: 33.1in
\setlength{\sepwid}{0.024\paperwidth} % Separation width (white space) between columns
\setlength{\onecolwid}{0.22\paperwidth} % Width of one column
\setlength{\twocolwid}{0.464\paperwidth} % Width of two columns
\setlength{\threecolwid}{0.708\paperwidth}
\setlength{\fourcolwid}{0.9508\paperwidth} % Width of three columns
\setlength{\topmargin}{-0.75in} % Reduce the top margin size
%-----------------------------------------------------------

\usepackage{graphicx}  % Required for including images
\usepackage[T1]{fontenc}
\usepackage{charter}
\usepackage[expert]{mathdesign}
\usepackage[percent]{overpic}
%\usepackage{caption}
\usepackage{booktabs} % Top and bottom rules for tables\usepackage[T1]{fontenc}
\usepackage{charter}

\DeclareFontShape{OMX}{cmex}{m}{n}{
  <-7.5> cmex7
  <7.5-8.5> cmex8
  <8.5-9.5> cmex9
  <9.5-> cmex10
}{}
\SetSymbolFont{largesymbols}{normal}{OMX}{cmex}{m}{n}
\SetSymbolFont{largesymbols}{bold}  {OMX}{cmex}{m}{n}


%----------------------------------------------------------------------------------------
%	TITLE SECTION 
%----------------------------------------------------------------------------------------

\title{\Large The Price of Fame:\\
\Large The Effect of Award-Winning Pedigrees on the Price of Livestock Genetics} % Poster title

\author{Victor Funes-Leal and Jared Hutchins} % Author(s)

\institute{University of Illinois at Urbana-Champaign\\ Department of Agricultural and Consumer Economics} % Institution(s)

%----------------------------------------------------------------------------------------

\begin{document}
\addtobeamertemplate{headline}{} 
{
%\begin{tikzpicture}[remember picture,overlay] 
%\node [shift={(-10 cm,-7cm)}] at (current page.north east) %{
\includegraphics[scale=2]{BlockI-Logo}}; 
%\end{tikzpicture} 

\begin{tikzpicture}[remember picture,overlay] 
	\node [shift={(14 cm,-7cm)}] at (current page.north west) {
\includegraphics[scale=3]{ACES-Logo}}; 
\end{tikzpicture} 
}

% {
% \begin{tikzpicture}[remember picture,overlay] 
% \node [shift={(-1 cm,-7cm)}] at (current page.north west) {\includegraphics[scale=.1]{bucky}}; 
% \end{tikzpicture} 
% }
\addtobeamertemplate{block end}{}{\vspace*{2ex}} % White space under blocks
\addtobeamertemplate{block alerted end}{}{\vspace*{2ex}} % White space under highlighted (alert) blocks

\setlength{\belowcaptionskip}{2ex} % White space under figures
\setlength\belowdisplayshortskip{2ex} % White space under equations

\begin{frame}[t] % The whole poster is enclosed in one beamer frame
\begin{columns}[t] % The whole poster consists of three major columns, the second of which is split into two columns twice - the [t] option aligns each column's content to the top

\begin{column}{\sepwid}\end{column} % Empty spacer column

\begin{column}{\onecolwid} % The first column

%----------------------------------------------------------------------------------------
%	OBJECTIVES
%----------------------------------------------------------------------------------------

\begin{block}{Introduction}
\begin{center}
%\includegraphics[scale=2]{milk_price.png} 
\end{center}
\begin{itemize}
\item \textbf{Research Question}: What is the causal effect of award-winning pedigrees on the price of livestock genetics?
\item \textbf{Natural experiment}: The price difference between World Dairy Expo winners and runner-ups is mostly random since the margin is relatively small.
\item Genetic traits are an important factor to determine a bull's semen price, particularly productivity and health traits (Schroeder et al., 1992; Richards and Jeffrey, 1996; Sy et , 1997).
\item None of these papers include pedigree into the pricing equation (focus on pricing each trait).
\end{itemize}

% \includegraphics[scale=2]{sample_milk_price.png}
\end{block}
\begin{block}{Data Description}
\begin{itemize}
\item 5367 quarterly evaluation observations from 476 distinct bulls belonging to five different lineages reconstructed from sire/dam data.
\item Ancestry data was merged with bull-level evaluation data from 2001 to 2004 (years of high winner turnover).
\item Source: National Association of Animal Breeders and winner/runner-up data from the World Dairy Expo (Held on October of every year).
\end{itemize}
\end{block}

%----------------------------------------------------------------------------------------
%	QUICK REVISION
%----------------------------------------------------------------------------------------

%\includegraphics[scale=1.75]{wood_model}

%----------------------------------------------------------------------------------------

\end{column} % End of the first column



\begin{column}{\sepwid}\end{column} % Empty spacer column

\begin{column}{\twocolwid} % Begin a column which is two columns wide (column 2)
	
\begin{columns}[t,totalwidth=\twocolwid]
	\begin{column}{\twocolwid}
		\begin{block}{Pedigree Networks}
			%\vspace{-0.9in}
			\begin{figure}
				\begin{center}
					\includegraphics[scale=1.9]{Network.pdf}\\
					\caption{\hspace{.05cm}WDE winners/runners-up and relatives}
					\label{fig:network}
				\end{center}
			\end{figure} 
		\end{block}
	\end{column}
\end{columns}
	

\begin{columns}[t,totalwidth=\twocolwid] % Split up the two columns wide column

\begin{column}{\onecolwid}\vspace{-.6in} % The first column within column 2 (column 2.1)

%----------------------------------------------------------------------------------------
%	MATERIALS
%----------------------------------------------------------------------------------------


%----------------------------------------------------------------------------------------
\begin{block}{Event Study 1}
	%\vspace{-0.9in}
	\begin{figure}
		\begin{center}
			\includegraphics[scale=1.2]{event_study.png}\\
			\caption{\hspace{.05cm}Weights proportional to genetic distance}
			\label{fig:study1}
		\end{center}
	\end{figure} 
\end{block}
\end{column} % End of column 2.1

\begin{column}{\onecolwid}\vspace{-.6in} % The second column within column 2 (column 2.2)
\begin{block}{Event Study 2}
		%\vspace{-0.9in}
		\begin{figure}
			\begin{center}
				\includegraphics[scale=1.2]{event_study_alt.png}\\
				\caption{\hspace{.05cm}Weights equal to 1 for winners and half-siblings only}
				\label{fig:study2}
			\end{center}
		\end{figure} 

\end{block}
\end{column}	
\end{columns} % End of the split of column 2 - any content after this will now take up 2 columns width
\begin{block}{ }
After the competition, the relatives of the winner increased \$2.75 in price, a 12\% increase, relative to the runners up. Half-siblings had a larger increase of \$3.14, though for all the relatives the price increase is temporary.
\end{block}
\end{column} % End of the second column

\begin{column}{\sepwid}\end{column} % Empty spacer column

\begin{column}{\onecolwid} % The third column
	
\begin{block}{Rotation Plans}

\begin{table}
    \centering
    \begin{tabular}{ccc}
        Code & Pattern  & \% Corn & Prevalence\\
       \textcolor{red}{CCCCCC}  & Corn monoculture & 100 & 53\%\\
        \twxtcolor{}{} &  & \\
    \end{tabular}
    \caption{Caption}
    \label{tab:placeholder}
\end{table}
\vspace{1.2cm}
\end{block}	

\begin{block}{Empirical Model}
\begin{equation*}
price_{it} = \sum_{k=-3,k\neq -1}^{3}\beta_{k}\tau_{t+k}w_{i}+\mathbf{X}_{it}'\Gamma+\gamma_{i}+\lambda_{t}+\epsilon_{it}
\end{equation*}

\quad \\

\vspace{-1cm}
\begin{itemize}
\item $price_{it}$: Semen price for bull $i$ at time $t$.
\item $\tau_{t+k}$: Lags and leads of the treatment indicator (1 if bull $i$ is part of the winner network at time $t=0$).
\item $w_{i}$: Genetic distance to the winner/runner-up, defined as $2^{-p}$ where $p=1$ for half-brothers, $p=2$ for cousins, etc. 
\item $\mathbf{X}_{it}$ is the set of genetic traits.
\item $\gamma_{i}$, $\lambda_{t}$: bull and quarter fixed effects.
\end{itemize}


\end{block}

\begin{block}{Conclusions}
\begin{itemize}
\item Reported bull semen price responds positively to winning the World Dairy Expo, bulls on the winner network have a 10\% increase in semen prices after the event.
\item Winning the WDE impacts close relatives as well, but its short-lived and disappears by the quarter previous to the next Expo.
\end{itemize}
\quad\\
\quad\\
\end{block}

\end{column} % End of the third column
\end{columns} % End of all the columns in the poster

\begin{columns}[t]
\begin{column}{\fourcolwid}
\begin{textblock}{4}(12.25,-.65)
\centering
 \textbf{Contact Info}\\
 Victor Funes: victorf2@illinois.edu\\
 Jared Hutchins: jhtchns2@illinois.edu\\
\end{textblock}
% \quad \\
%\vspace{-2cm}
\footnotesize
\textbf{References}\\
\footnotesize
Schroeder, T., Espinosa, J. and Goodwin, B. (1992). The Value of Genetic Traits in Purebred Dairy Bull Services. \emph{Review of Agricultural Economics 14} (2), 215--226.\\
Richards, T. and Jeffrey, S. (1996) Establishing Indices of genetic Merit Using Hedonic Pricing: An Application to dairy Bulls in Alberta. \emph{Canadian Journal of Agricultural Economics 44}, 251--264. \\
Sy, H., Faminow, M. Johnson, G. and Crow, G. (1997). Estimating the Values of Cattle Characteristics Using an Ordered Probit Model. \emph{American Journal of Agricultural Economics 79} (2), 463--476.\\
\end{column}
\end{columns}
\end{frame} % End of the enclosing frame
\end{document}
